% Section 2, part 2 of 8: the large-radius cyclic structure (statement; grading, unit, periods, Birkhoff recursion).

Along the line $t=T(q)H$, with $\Qn=1$, functions of $t$ are functions of $q$, and $\partial_T\coloneqq C_1^{-1}\Dq$ is the derivative in the coordinate $T$ of $H$, i.e.\ $\partial_{t^1}$ along the line ($\partial_T$ is not the derivation $\de_t$ of \cref{def:gau-conc}, which is defined in a conifold frame). By (G0), the quantum product $*$ at $t=TH$ is commutative and associative, has the unit $1$, and is Frobenius for $\eta$: $\eta(x*y,w)=\langle\!\langle x,y,w\rangle\!\rangle=\eta(x,y*w)$.

\begin{prop}[The large-radius cyclic structure]
\label{prop:set-cyclicLR}
Assume (M) and (G0). Along the line $t=T(q)H$, at $\lambda_i=\zeta^i$ and $\Qn=1$, the following hold.
\begin{enumerate}
\item In the flat basis, $H*H^a=\kappa^{\mathrm{LR}}_aH^{a+1}$ for $a\in\bbZ/5$, where
\[
\kappa^{\mathrm{LR}}=(\kappa^{\mathrm{LR}}_0,\dots,\kappa^{\mathrm{LR}}_4)=\Bigl(1,\ \frac{L^5}{C_0^2C_1^3},\ 1,\ \frac{C_0}{C_1},\ \frac{C_0}{C_1}\Bigr).
\]
\item Consequently, for every solution $v$ of the quantum differential equation $z\partial_Tv=H*v$, in any differential $\bbQ\fps q$-algebra in which $z$ is a constant, the lowered components $v_a\coloneqq\eta(H^a,v)$ satisfy $z\Dq v_a=C_1\kappa^{\mathrm{LR}}_a\,v_{a+1}$ for all $a\in\bbZ/5$.
\item For every colour $i$ the rows of \cref{def:set-rows} (defined in \cref{sec:set-rows}) satisfy $(z\Dq+\zeta^iL)\Qr^{(i)}_a=C_1\kappa^{\mathrm{LR}}_a\,\Qr^{(i)}_{a+1}$ for $a\in\bbZ/5$. Equivalently, for $X_a\coloneqq\zeta^{-i}\Qr^{(i)}_a$,
\[
X_0=L\Phi_i/C_0,\qquad X_{a+1}=(C_1\kappa^{\mathrm{LR}}_a)^{-1}(z\Dq+\zeta^iL)X_a\quad(a=0,\dots,3),\qquad(z\Dq+\zeta^iL)X_4=C_0X_0 .
\]
\end{enumerate}
\end{prop}

The superscript LR distinguishes $\kappa^{\mathrm{LR}}$ from the coefficients $\kappa_a$ of the cyclic system in a conifold frame (\cref{prop:gau-quintic}). The proof of (i) and (ii) occupies the rest of this subsection up to \cref{cor:set-cyclicPF5}; it uses (M), (G0), \cref{lem:set-pf5} and the pairing of \cref{def:set-fq}, through six lemmas. Part (iii) concerns the rows, which are defined only after \cref{cor:set-cyclicPF5}~(a); it is proved after \cref{def:set-rows}, from (ii) and \cref{cor:set-cyclicPF5}~(a). The letters $\mathfrak w$, $\cP$, $\mathfrak p_k$, $\tilde{\mathfrak p}_m$, $\mathsf p^{(a)}_k$, $X_a$, $\Mcyc^{\mathrm{LR}}$ and $\chi$ are local to this subsection. Before \cref{prop:set-cyclicLR}~(i) is proved, $\kappa^{\mathrm{LR}}_a$ denotes the structure constant of \cref{lem:set-grading}.

\begin{lem}[The $\bbZ/5$ grading]
\label{lem:set-grading}
At $\lambda_i=\zeta^i$ and $t=TH$, $\langle\!\langle H^a,H^b,H^c\rangle\!\rangle=0$ unless $a+b+c\equiv3\bmod5$. Hence $H^a*H^b$ is a multiple of $H^{a+b}$, and $H*H^a=\kappa^{\mathrm{LR}}_aH^{a+1}$ with $\kappa^{\mathrm{LR}}_a\in1+q\bbQ\fps q$.
\end{lem}

\begin{proof}
By the divisor equation (G0)(b), with $\Qn=1$,
\[
\langle\!\langle H^a,H^b,H^c\rangle\!\rangle(TH)=\eta(H^aH^b,H^c)+\sum_{d\geq1}\re^{dT}\langle H^a,H^b,H^c\rangle_{0,3,d};
\]
in degree $0$ the insertions of $TH$ contribute nothing, because $\overline{\mathcal M}_{0,3+n}(\bbP^4,0)=\overline{\mathcal M}_{0,3+n}\times\bbP^4$ and the integrand is pulled back from $\bbP^4$. By (G0)(a), before the specialisation the term of degree $d\geq1$ is a symmetric polynomial in $\lambda_0,\dots,\lambda_4$, homogeneous of degree $m=a+b+c-3$, that is, a polynomial in the elementary symmetric functions $e_1,\dots,e_5$, of weighted degree $m$. At $\lambda_j=\zeta^j$ the identity $\prod_j(x-\zeta^j)=x^5-1$ gives $e_1=\dots=e_4=0$ and $e_5=1$, so the value is $0$ unless $5\mid m$, i.e.\ unless $a+b+c\equiv3\bmod5$. The classical term is $\eta(H^{a+b},H^c)=5\,\delta_{a+b+c\equiv3}$ (\cref{def:set-fq}). Since $\eta^{cc'}$ is supported on $c+c'\equiv3$, the product $H^a*H^b=\sum_{c,c'}\langle\!\langle H^a,H^b,H^c\rangle\!\rangle\eta^{cc'}H^{c'}$ is supported on $c'\equiv a+b$. The structure constant $\kappa^{\mathrm{LR}}_a$ is a power series in $\re^T$ with rational coefficients, hence lies in $\bbQ\fps q$ through $\re^T=q\exp(\tilde I_1/C_0)$; at $q=0$ the product is the classical one, $H*H^a=H^{a+1}$, so $\kappa^{\mathrm{LR}}_a(0)=1$.
\end{proof}

\begin{lem}[Unit and Frobenius property]
\label{lem:set-unitfrob}
$\kappa^{\mathrm{LR}}_0=\kappa^{\mathrm{LR}}_2=1$ and $\kappa^{\mathrm{LR}}_3=\kappa^{\mathrm{LR}}_4$.
\end{lem}

\begin{proof}
The class $1$ is the unit of $*$ by the string equation, so $H*1=H$ and $\kappa^{\mathrm{LR}}_0=1$. By commutativity and the Frobenius property, $\eta(H*H^a,H^{2-a})=\eta(H^a,H*H^{2-a})$, i.e.\ $5\kappa^{\mathrm{LR}}_a=\eta(\kappa^{\mathrm{LR}}_aH^{a+1},H^{2-a})=\eta(H^a,\kappa^{\mathrm{LR}}_{2-a}H^{3-a})=5\kappa^{\mathrm{LR}}_{2-a}$ (\cref{def:set-fq}). So $\kappa^{\mathrm{LR}}_a=\kappa^{\mathrm{LR}}_{2-a}$ for all $a\in\bbZ/5$, that is $\kappa^{\mathrm{LR}}_0=\kappa^{\mathrm{LR}}_2$ and $\kappa^{\mathrm{LR}}_3=\kappa^{\mathrm{LR}}_4$.
\end{proof}

\begin{lem}[The periods]
\label{lem:set-periods}
Let $\mathfrak w$ be an indeterminate and
\[
\cP(\mathfrak w,q)\coloneqq\sum_{d\geq0}q^d\,\frac{\prod_{k=1}^{5d}(5\mathfrak w+k)}{\prod_{k=1}^d\bigl((\mathfrak w+k)^5-\mathfrak w^5\bigr)}\in\bbQ\fps q\fps{\mathfrak w},
\]
the denominators being polynomials in $\mathfrak w$ with the constant term $(d!)^5$. The map $\bbQ\fps q[\logq]\fps{\mathfrak w}\to\bbH\otimes\bbQ\fps q[\logq]\fps{z^{-1}}$, $\sum_ka_k\mathfrak w^k\mapsto\sum_ka_kH^kz^{-k}$ (with $H^k\coloneqq H^{k\bmod5}$), is an injective ring homomorphism, and $I/z$ is the image of $\re^{\logq\mathfrak w}\cP(\mathfrak w,q)$. Hence
\[
\frac Iz=\sum_{k\geq0}\mathfrak p_kH^kz^{-k},\qquad\mathfrak p_k\coloneqq\sum_{j=0}^k\frac{\logq^j}{j!}\,\tilde{\mathfrak p}_{k-j},\qquad\tilde{\mathfrak p}_m\coloneqq\coef{\mathfrak w^m}\cP\in\bbQ\fps q,\quad\tilde{\mathfrak p}_m(0)=\delta_{m0},
\]
with $\mathfrak p_0=C_0$ and $\mathfrak p_1=I_1$. Moreover $\PFfour\mathfrak p_k=0$ for $k=0,\dots,3$, and $\PFfour\mathfrak p_4=1$.
\end{lem}

\begin{proof}
The map is multiplicative because $H^5=1$, and injective because distinct powers of $\mathfrak w$ go to distinct powers of $z$. Followed by the restriction to $p_i$ it sends a series $G(\mathfrak w)$ to $G(\zeta^i/z)$. At the zeta point $\prod_j(\zeta^i-\zeta^j+kz)=(\zeta^i+kz)^5-1$, and $\prod_{k=1}^{5d}(5\zeta^i/z+k)=z^{-5d}\prod_{k=1}^{5d}(5\zeta^i+kz)$, $\prod_{k=1}^d((\zeta^i/z+k)^5-(\zeta^i/z)^5)=z^{-5d}\prod_{k=1}^d((\zeta^i+kz)^5-1)$; so the powers of $z$ cancel in each term $q^d$, and as $q^{\zeta^i/z}=\re^{\logq\zeta^i/z}$, the restriction of the image of $\re^{\logq\mathfrak w}\cP$ to $p_i$ is $I_i/z$. As the five restrictions determine an element, $I/z$ is the image of $\re^{\logq\mathfrak w}\cP$, and expanding $\re^{\logq\mathfrak w}\cP=\sum_k\mathfrak p_k\mathfrak w^k$ gives the formula for $\mathfrak p_k$. The term $q^0$ of $\cP$ is $1$, so $\tilde{\mathfrak p}_m(0)=\delta_{m0}$. The constant term of $\cP$ in $\mathfrak w$ is $\sum_dq^d(5d)!/(d!)^5=C_0$, and its coefficient of $\mathfrak w$ is, by logarithmic differentiation, $\sum_dq^d\frac{(5d)!}{(d!)^5}\bigl(\sum_{k=1}^{5d}\frac5k-\sum_{k=1}^d\frac5k\bigr)=\tilde I_1$; hence $\mathfrak p_0=C_0$ and $\mathfrak p_1=C_0\logq+\tilde I_1=I_1$.

\emph{The operator $\PFfive$.} Since $z$ is a constant, $\prod_{j}(z\Dq-\zeta^j)=(z\Dq)^5-1$. On the term $q^d$ of $I_i$, which is $q^{d+\zeta^i/z}$ times a function of $z$, the operator $z\Dq$ acts by $\zeta^i+dz$, so $\prod_j(z\Dq-\zeta^j)$ multiplies it by $\prod_j(\zeta^i-\zeta^j+dz)$; this cancels the factors $k=d$ of the denominator (and kills the term $d=0$), and leaves $q\prod_{k=1}^5(5\zeta^i+(5d-5+k)z)$ times the term $q^{d-1}$. On the term $q^{d-1}$ the operator $\prod_{k=1}^5(5z\Dq+kz)$ acts by the same factor $\prod_{k=1}^5(5\zeta^i+(5d-5+k)z)$. Hence $\PFfive I_i=0$. By \cref{lem:set-pf5} and $I_i/z=\sum_k\zeta^{ik}\mathfrak p_kz^{-k}$,
\[
0=\PFfive(I_i/z)=\sum_{k\geq0}\zeta^{ik}\bigl(z^{5-k}\,\Dq\PFfour\mathfrak p_k-z^{-k}\mathfrak p_k\bigr),
\]
and the coefficients of $z^5,z^4,\dots,z^1$ give $\Dq\PFfour\mathfrak p_k=0$ for $k=0,\dots,4$.

\emph{Constants.} The constants of $\Dq$ on $\bbQ\fps q[\logq]$ are $\bbQ$: on $\bbQ\fps q$, $\Dq\sum_nc_nq^n=\sum_nnc_nq^n$, whose constants are $\bbQ$; and if $f=\sum_{j\leq N}a_j\logq^j$ with $a_j\in\bbQ\fps q$, $a_N\neq0$, $N\geq1$ and $\Dq f=0$, the coefficient of $\logq^N$ gives $\Dq a_N=0$, so $a_N\in\bbQ$, and the coefficient of $\logq^{N-1}$ gives $\Dq a_{N-1}=-Na_N\neq0$, which is impossible since $\Dq$ maps $\bbQ\fps q$ into $q\bbQ\fps q$. Hence $\PFfour\mathfrak p_k\in\bbQ$ for $k\leq4$, equal to its part of order $q^0$. Taking the part of order $q^0$ commutes with $\Dq$, which acts on it as $\dd/\dd\logq$, and it kills the multiples of $q$. By \eqref{eq:set-PF4} and $\tilde{\mathfrak p}_m(0)=\delta_{m0}$, the part of order $q^0$ of $\PFfour\mathfrak p_k$ is $(\dd/\dd\logq)^4(\logq^k/k!)$, which is $0$ for $k\leq3$ and $1$ for $k=4$.
\end{proof}

\begin{lem}[The Birkhoff recursion]
\label{lem:set-birkhoff}
Put $\mathsf p^{(0)}_k\coloneqq\mathfrak p_k/C_0$ for $k\geq0$ and, recursively for $a=0,\dots,3$, $\mathsf p^{(a+1)}_k\coloneqq\partial_T\mathsf p^{(a)}_k/\kappa^{\mathrm{LR}}_a$ for $k\geq a+1$. Then, along the line $t=TH$,
\[
J_a=\sum_{k\geq a}\mathsf p^{(a)}_k\,H^kz^{a-k},\qquad\mathsf p^{(a)}_a=1\qquad(a=0,\dots,4),
\]
and $\kappa^{\mathrm{LR}}_a=\partial_T\mathsf p^{(a)}_{a+1}$ for $a=0,\dots,3$. In particular $\mathsf p^{(0)}_1=T$, $\kappa^{\mathrm{LR}}_1=\partial_T^2\mathsf p^{(0)}_2$, $\kappa^{\mathrm{LR}}_2=\partial_T\bigl(\partial_T^2\mathsf p^{(0)}_3/\kappa^{\mathrm{LR}}_1\bigr)$ and, since $\kappa^{\mathrm{LR}}_2=1$, $\kappa^{\mathrm{LR}}_3=\partial_T^2\bigl(\partial_T^2\mathsf p^{(0)}_4/\kappa^{\mathrm{LR}}_1\bigr)$.
\end{lem}

\begin{proof}
Induction on $a$. By (G0)(d) and (M), $J_0=J/z=I/(zC_0)$ along the line, and by \cref{lem:set-periods} this is $\sum_k(\mathfrak p_k/C_0)H^kz^{-k}$, with $\mathsf p^{(0)}_0=\mathfrak p_0/C_0=1$; distinct $k$ carry distinct powers of $z$, so the coefficients are well defined. By (G0)(d) with $b=1$ and \cref{lem:set-grading}, $z\partial_TJ_a=\kappa^{\mathrm{LR}}_aJ_{a+1}$ for $a=0,\dots,3$, the derivative in the direction $H$ along the line being $\partial_T$. Differentiating the expansion of $J_a$ gives $z\partial_TJ_a=\sum_{k\geq a+1}(\partial_T\mathsf p^{(a)}_k)H^kz^{a+1-k}$, the term $k=a$ being $z\partial_T1=0$. Since $J_{a+1}=H^{a+1}+O(z^{-1})$, the coefficient of $z^0$ gives $\kappa^{\mathrm{LR}}_a=\partial_T\mathsf p^{(a)}_{a+1}$; as $\kappa^{\mathrm{LR}}_a$ is a unit of $\bbQ\fps q$ (\cref{lem:set-grading}), $J_{a+1}=z\partial_TJ_a/\kappa^{\mathrm{LR}}_a$ gives the expansion of $J_{a+1}$, with $\mathsf p^{(a+1)}_{a+1}=1$. For the closed forms, $\mathsf p^{(0)}_1=I_1/C_0=T$, and $\kappa^{\mathrm{LR}}_0=\partial_TT=1$ gives $\mathsf p^{(1)}_k=\partial_T\mathsf p^{(0)}_k$; then $\mathsf p^{(2)}_k=\partial_T^2\mathsf p^{(0)}_k/\kappa^{\mathrm{LR}}_1$ and $\mathsf p^{(3)}_k=\partial_T\mathsf p^{(2)}_k/\kappa^{\mathrm{LR}}_2$, and the formulas follow from $\kappa^{\mathrm{LR}}_a=\partial_T\mathsf p^{(a)}_{a+1}$ and \cref{lem:set-unitfrob}.

Only derivatives along the line $t=TH$ are used: $J_{a+1}$ is obtained from $J_a$ by the quantum differential equation in the direction $b=1$, so the J-function is never needed off the line.
\end{proof}

% Section 2, part 3 of 8: the Wronskian, operators of order three, proof of prop:set-cyclicLR, its corollary; points; Frobenius data.

\begin{lem}[The Wronskian]
\label{lem:set-wronskian}
$\Wr(\mathfrak p_0,\dots,\mathfrak p_3)\coloneqq\det\bigl(\Dq^j\mathfrak p_k\bigr)_{j,k=0,\dots,3}=L^{10}$.
\end{lem}

\begin{proof}
By \eqref{eq:set-PF4} the leading coefficient of $\PFfour$ is $1-5^5q=L^{-5}$, a unit, and its coefficient of $\Dq^3$ is $-6250q$. Differentiate the determinant row by row: the rows $j<3$ give determinants with two equal rows, and in the last row $\Dq^4\mathfrak p_k=-L^5\bigl(-6250q\,\Dq^3\mathfrak p_k+(\text{terms in }\Dq^2\mathfrak p_k,\Dq\mathfrak p_k,\mathfrak p_k)\bigr)$ for $k\leq3$ (\cref{lem:set-periods}), whose last terms are combinations of the other rows. Hence $\Dq\Wr=6250qL^5\,\Wr$. Also $\Dq L^{10}=10L^9\Dq L=6250qL^{15}=6250qL^5\,L^{10}$, by $\Dq L=5^4qL^6$. So $\Dq(\Wr\,L^{-10})=0$, and $\Wr\,L^{-10}$ is a constant (the constants of $\Dq$ on $\bbQ\fps q[\logq]$ are $\bbQ$, proof of \cref{lem:set-periods}), equal to its part of order $q^0$. Taking this part is a ring homomorphism to $\bbQ[\logq]$ on which $\Dq$ acts as $\dd/\dd\logq$; by $\tilde{\mathfrak p}_m(0)=\delta_{m0}$ it maps $\Dq^j\mathfrak p_k$ to $\logq^{k-j}/(k-j)!$ for $k\geq j$ and to $0$ for $k<j$, so it maps the determinant to that of an upper unitriangular matrix, which is $1$.
\end{proof}

\begin{lem}[Operators of order at most three]
\label{lem:set-order3}
If $\Theta=\sum_{j=0}^3r_j\Dq^j$ with $r_j\in\bbQ\lser q$ satisfies $\Theta\mathfrak p_k=0$ for $k=0,\dots,3$, then $\Theta=0$.
\end{lem}

\begin{proof}
By \cref{lem:set-periods}, $\sum_k\mathfrak p_k\mathfrak w^k=\re^{\logq\mathfrak w}\cP(\mathfrak w,q)$. Since the $r_j$ do not involve $\logq$ and $\Dq\re^{\logq\mathfrak w}=\mathfrak w\re^{\logq\mathfrak w}$, the Leibniz rule gives, coefficientwise in $\mathfrak w$,
\[
\sum_k(\Theta\mathfrak p_k)\mathfrak w^k=\Theta\bigl(\re^{\logq\mathfrak w}\cP\bigr)=\re^{\logq\mathfrak w}\,\cP_\Theta,\qquad\cP_\Theta\coloneqq\sum_{j=0}^3r_j(\Dq+\mathfrak w)^j\cP\in\bbQ\lser q\fps{\mathfrak w}.
\]
The hypothesis says that the left side is $O(\mathfrak w^4)$; as $\re^{\logq\mathfrak w}$ is invertible with constant term $1$, $\cP_\Theta=O(\mathfrak w^4)$. Now $\coef{\mathfrak w^n}\cP_\Theta=\sum_{j=0}^3m_{nj}r_j$ with $m_{nj}\coloneqq\sum_{i\leq\min(j,n)}\binom ji\Dq^{j-i}\tilde{\mathfrak p}_{n-i}\in\bbQ\fps q$, $n,j=0,\dots,3$. Modulo $q$, $\Dq^e\tilde{\mathfrak p}_m\equiv0$ unless $e=0$ and $m=0$, because $\tilde{\mathfrak p}_m(0)=\delta_{m0}$ and $\Dq$ kills $q^0$; so $m_{nj}\equiv\delta_{nj}\bmod q$. Hence $\det(m_{nj})$ is a unit of $\bbQ\fps q$, and $\sum_jm_{nj}r_j=0$ for $n=0,\dots,3$ gives $r_j=0$.
\end{proof}

\begin{proof}[Proof of \cref{prop:set-cyclicLR}]
(i) By \cref{lem:set-unitfrob}, $\kappa^{\mathrm{LR}}_0=\kappa^{\mathrm{LR}}_2=1$ and $\kappa^{\mathrm{LR}}_3=\kappa^{\mathrm{LR}}_4$.

\emph{The coefficient $\kappa^{\mathrm{LR}}_1$.} By \cref{lem:set-periods,lem:set-birkhoff}, $\mathfrak p_k=C_0\mathsf p^{(0)}_k$ with $\mathsf p^{(0)}_0=1$ and $\mathsf p^{(0)}_1=T$. For any four functions $f_k$, $\det(\Dq^j(C_0f_k))=C_0^4\det(\Dq^jf_k)$, because $\Dq^j(C_0f)$ is $C_0\Dq^jf$ plus a combination of the $\Dq^{j'}f$, $j'<j$, with coefficients independent of $f$; and $\det(\Dq^jf_k)=C_1^{0+1+2+3}\det(\partial_T^jf_k)$, because $\Dq=C_1\partial_T$, so that $\Dq^j$ is $C_1^j\partial_T^j$ plus terms with fewer $\partial_T$. With $f=(1,T,\mathsf p^{(0)}_2,\mathsf p^{(0)}_3)$ and $\partial_TT=1$,
\[
\det(\partial_T^jf_k)=\partial_T^2\mathsf p^{(0)}_2\,\partial_T^3\mathsf p^{(0)}_3-\partial_T^3\mathsf p^{(0)}_2\,\partial_T^2\mathsf p^{(0)}_3=\bigl(\partial_T^2\mathsf p^{(0)}_2\bigr)^2\,\partial_T\Bigl(\frac{\partial_T^2\mathsf p^{(0)}_3}{\partial_T^2\mathsf p^{(0)}_2}\Bigr)=(\kappa^{\mathrm{LR}}_1)^2\kappa^{\mathrm{LR}}_2
\]
by \cref{lem:set-birkhoff}. So \cref{lem:set-wronskian} gives $C_0^4C_1^6(\kappa^{\mathrm{LR}}_1)^2\kappa^{\mathrm{LR}}_2=L^{10}$, and with $\kappa^{\mathrm{LR}}_2=1$,
\[
\Bigl(\kappa^{\mathrm{LR}}_1-\frac{L^5}{C_0^2C_1^3}\Bigr)\Bigl(\kappa^{\mathrm{LR}}_1+\frac{L^5}{C_0^2C_1^3}\Bigr)=0
\]
in the domain $\bbQ\fps q$. The second factor has the constant term $2$ (\cref{lem:set-grading}; $C_0(0)=C_1(0)=L(0)=1$), so $\kappa^{\mathrm{LR}}_1=L^5/(C_0^2C_1^3)$.

\emph{The coefficient $\kappa^{\mathrm{LR}}_3$.} Let $\Mcyc^{\mathrm{LR}}\coloneqq\partial_T^2\circ(\kappa^{\mathrm{LR}}_1)^{-1}\circ\partial_T^2\circ C_0^{-1}$, a differential operator with coefficients in $\bbQ\fps q$ ($C_0$, $C_1$ and $\kappa^{\mathrm{LR}}_1$ are units). By \cref{lem:set-birkhoff}, $\Mcyc^{\mathrm{LR}}\mathfrak p_0=\partial_T^2(\kappa^{\mathrm{LR}}_1)^{-1}\partial_T^21=0$, $\Mcyc^{\mathrm{LR}}\mathfrak p_1=0$ (as $\partial_T^2T=0$), $\Mcyc^{\mathrm{LR}}\mathfrak p_2=\partial_T^2\bigl((\kappa^{\mathrm{LR}}_1)^{-1}\kappa^{\mathrm{LR}}_1\bigr)=0$, $\Mcyc^{\mathrm{LR}}\mathfrak p_3=\partial_T\kappa^{\mathrm{LR}}_2=0$ and $\Mcyc^{\mathrm{LR}}\mathfrak p_4=\kappa^{\mathrm{LR}}_3$. In powers of $\Dq$, $\Mcyc^{\mathrm{LR}}=(\kappa^{\mathrm{LR}}_1C_0C_1^4)^{-1}\Dq^4+(\text{order}\leq3)$, and $\PFfour=L^{-5}\Dq^4+(\text{order}\leq3)$ by \eqref{eq:set-PF4}. So $\Mcyc^{\mathrm{LR}}-\chi\,\PFfour$, $\chi\coloneqq L^5/(\kappa^{\mathrm{LR}}_1C_0C_1^4)$, has order at most $3$ and kills $\mathfrak p_0,\dots,\mathfrak p_3$ (\cref{lem:set-periods}); by \cref{lem:set-order3} it vanishes:
\[
\Mcyc^{\mathrm{LR}}=\chi\,\PFfour .
\]
Hence $\kappa^{\mathrm{LR}}_3=\Mcyc^{\mathrm{LR}}\mathfrak p_4=\chi\,\PFfour\mathfrak p_4=\chi$ (\cref{lem:set-periods}), and by the value of $\kappa^{\mathrm{LR}}_1$, $\chi=C_0^2C_1^3/(C_0C_1^4)=C_0/C_1$. Finally $\kappa^{\mathrm{LR}}_4=\kappa^{\mathrm{LR}}_3$. This proves (i).

(ii) For a solution $v$ of the quantum differential equation the Frobenius property and the commutativity of $*$ give $z\partial_T\eta(H^a,v)=\eta(H^a,H*v)=\eta(H*H^a,v)=\kappa^{\mathrm{LR}}_a\eta(H^{a+1},v)$, and $z\partial_T=C_1^{-1}z\Dq$; the pairing $\eta(H^a,\cdot)$ has constant coefficients, so it commutes with $\partial_T$.

(iii) is proved after \cref{def:set-rows}, which defines the rows it concerns.
\end{proof}

Every step is an identity in $\bbQ\fps q$ or in $\bbQ\fps q[\logq]$; no truncation in $q$ and no analytic continuation is used. The proof uses (M), (G0)(a)--(d), \cref{lem:set-pf5} and the restrictions of $I$ to the fixed points (in the derivation of $\Dq\PFfour\mathfrak p_k=0$ in \cref{lem:set-periods} one colour suffices); it uses neither the cyclic form (CYC) of \cref{thm:gau-cyc} nor any statement in a conifold frame. The operator $\Mcyc^{\mathrm{LR}}$ is the large-radius analogue of the operator $\Mcyc$ of \cref{def:gau-conc}. The coefficient $\kappa^{\mathrm{LR}}_3=\Mcyc^{\mathrm{LR}}\mathfrak p_4$ is the one that sees the equivariant relation $H^5=1$: it involves the period $\mathfrak p_4$, which satisfies the inhomogeneous equation $\PFfour\mathfrak p_4=1$. Relations equivalent to (i) are stated, in other notation, in \cite[Props.~6.1--6.2]{GR19} and, for $a\leq3$, in \cite[Section~2.4.3, (17), and Section~2.4.4, (19)--(20)]{LP20}; the proof above does not use them.

\begin{cor}[Eigenvalues, idempotents and the fifth-order equation]
\label{cor:set-cyclicPF5}
Assume (M) and (G0).
\begin{itemize}
\item[(a)] $(H*)^5=(L/C_1)^5$. The eigenvalues of $H*$ are $\mu_i\coloneqq\zeta^iL/C_1$ ($i\in\bbZ/5$), labelled by $\mu_i(0)=\zeta^i$, the classical eigenvalue at $p_i$. The idempotents are
\[
e_i=\frac15\sum_{a=0}^4\Bigl(\prod_{b<a}\kappa^{\mathrm{LR}}_b\Bigr)\mu_i^{-a}H^a,\qquad\Delta_i\coloneqq\frac1{\eta(e_i,e_i)}=\frac{C_0^2\,\zeta^{3i}}{L^2}.
\]
The canonical coordinates along the line are functions $u_i$ with $\partial_Tu_i=\mu_i$, that is $\Dq u_i=\zeta^iL$; they are determined up to additive constants, the constants of $\Dq$ on $\bbQ(\zeta)\fps q[\logq]$ being $\bbQ(\zeta)$ (as in the proof of \cref{lem:set-periods}; for the conifold chart see \cref{lem:set-canon}).
\item[(b)] For every solution $v$ of $z\partial_Tv=H*v$ in a differential $\bbQ\fps q$-algebra in which $z$ is a constant, $C_0v_0$ solves $\PFfive$. In particular, for the column $v=R(z)\tilde e_i\re^{u_i/z}$ of \cref{def:set-rows}, the function $\re^{u_i/z}\,z\,L\,\Phi_i(\zh,u)$ solves $\PFfive$: this is the WKB form of colour $i$.
\end{itemize}
\end{cor}

\begin{proof}
(a) By \cref{prop:set-cyclicLR}~(i), $(H*)^5H^a=\kappa^{\mathrm{LR}}_a\kappa^{\mathrm{LR}}_{a+1}\cdots\kappa^{\mathrm{LR}}_{a+4}H^a$ and $\prod_a\kappa^{\mathrm{LR}}_a=L^5/C_1^5$. So $H*$ is the weighted cyclic shift with characteristic polynomial $x^5-(L/C_1)^5$; since $L/C_1\in1+q\bbQ\fps q$, its roots are the $\zeta^iL/C_1$, pairwise distinct. The eigenvector equation $H*x=\mu_ix$ for $x=\sum_ax_aH^a$ reads $x_a=x_{a-1}\kappa^{\mathrm{LR}}_{a-1}/\mu_i$, so $x_a=x_0\bigl(\prod_{b<a}\kappa^{\mathrm{LR}}_b\bigr)\mu_i^{-a}$, consistently around the cycle because $\prod_b\kappa^{\mathrm{LR}}_b=\mu_i^5$. As $*$ is commutative and associative with unit $1$, each eigenline of $H*$ is an ideal, and the five eigenlines, the eigenvalues being distinct, are spanned by the idempotents $e_i$, with $\sum_ie_i=1$. Write $x_0^{(i)}$ for the coefficient $x_0$ of $e_i$. The coefficient of $H^a$ in $\sum_ie_i=1$ is $\bigl(\prod_{b<a}\kappa^{\mathrm{LR}}_b\bigr)(C_1/L)^a\sum_ix_0^{(i)}\zeta^{-ia}=\delta_{a0}$ for $0\leq a\leq4$; as the matrix $(\zeta^{-ia})_{a,i}$ is invertible and $\sum_i\zeta^{-ia}=5\,\delta_{a0}$, the unique solution is $x_0^{(i)}=1/5$ for every $i$. By the Frobenius property $\eta(e_i,e_i)=\eta(1,e_i*e_i)=\eta(1,e_i)$, and $\eta(1,H^a)=5\,\delta_{a3}$ (\cref{def:set-fq}) gives $\eta(1,e_i)=\kappa^{\mathrm{LR}}_0\kappa^{\mathrm{LR}}_1\kappa^{\mathrm{LR}}_2\,\mu_i^{-3}=\kappa^{\mathrm{LR}}_1\mu_i^{-3}=L^2/(C_0^2\zeta^{3i})$. The canonical coordinates satisfy $\partial_Tu_i=\mu_i$, the eigenvalue of $H*$ on $e_i$; as $\Dq=C_1\partial_T$, this is $\Dq u_i=\zeta^iL$.

(b) From $z\partial_Tv_a=\kappa^{\mathrm{LR}}_av_{a+1}$ for $a=0,\dots,3$ (\cref{prop:set-cyclicLR}~(ii)) one gets $v_1=z\partial_Tv_0$, $v_2=z^2(\kappa^{\mathrm{LR}}_1)^{-1}\partial_T^2v_0$, $v_3=z^3\partial_T\bigl((\kappa^{\mathrm{LR}}_1)^{-1}\partial_T^2v_0\bigr)$ and $v_4=z^4(\kappa^{\mathrm{LR}}_3)^{-1}\Mcyc^{\mathrm{LR}}(C_0v_0)$. In the proof of \cref{prop:set-cyclicLR}, $\Mcyc^{\mathrm{LR}}=\chi\,\PFfour$ as operators and $\chi=\kappa^{\mathrm{LR}}_3$, so $v_4=z^4\,\PFfour(C_0v_0)$. The closing relation $z\partial_Tv_4=\kappa^{\mathrm{LR}}_4v_0$ becomes $z^5C_1^{-1}\Dq\PFfour(C_0v_0)=(C_0/C_1)v_0$, i.e.\ $(z^5\Dq\PFfour-1)(C_0v_0)=0$, which is $\PFfive(C_0v_0)=0$ by \cref{lem:set-pf5}. The last statement concerns the rows and is proved after \cref{def:set-rows}.
\end{proof}

The statement that the rows lie in the differential ring of the periods, the third consequence of \cref{prop:set-cyclicLR}, is \cref{cor:gau-rowsNN}.

\begin{defn}[Points and variables]
\label{def:set-points}
Put $u\coloneqq1-5^5q$ and $s\coloneqq u^{1/5}$, so that $L=u^{-1/5}=s^{-1}$ and $\Dq=q\,\dd/\dd q=-(1-u)\,\dd/\dd u$, and put $\zh\coloneqq z/s^4=L^4z$, the hatted variable. The \emph{large-radius point} is $q=0$ ($u=1$, $L=s=1$, $\zh=z$), the \emph{conifold point} is $u=0$ ($q=5^{-5}$), and the \emph{orbifold point} is $q=\infty$ ($L=0$). Near $q=0$ the functions of the note are elements of $\bbQ\fps q$ or $\bbQ\fps q[\logq]$; near the conifold point they are Laurent series in $s$, with $u=s^5$, $L=s^{-1}$ and $\Dq=-\frac{1-s^5}5s^{-4}\frac{\dd}{\dd s}$. The passage from the first description to the second is the conifold expansion of \cref{def:set-confexp}.
\end{defn}

\begin{defn}[Frobenius data and the topological field theory]
\label{def:set-frob}
At $\lambda_i=\zeta^i$ the pairing is (\cref{def:set-fq})
\[
\eta(H^a,H^b)=\sum_{i\in\bbZ/5}\frac{5\,\zeta^{i(a+b+1)}}{\prod_{j\neq i}(\zeta^i-\zeta^j)}=5\,\delta_{a+b\equiv3\bmod5},
\]
since $\prod_{j\neq i}(\zeta^i-\zeta^j)=5\zeta^{4i}$ is the derivative of $x^5-1$ at $\zeta^i$; as matrices $\eta^{-1}=\eta/25$, i.e.\ $\eta^{cc'}=\frac15\delta_{c+c'\equiv3}$. Along the line $t=T(q)H$ let $e_i$ be the idempotents and $\Delta_i=1/\eta(e_i,e_i)$ (\cref{cor:set-cyclicPF5}~(a)). The \emph{normalised canonical basis} is $\tilde e_i\coloneqq\Delta_i^{1/2}e_i$, so $\eta(\tilde e_i,\tilde e_j)=\delta_{ij}$, with the square roots
\[
\zeta^{1/2}\coloneqq\zeta^3,\qquad\Delta_i^{1/2}\coloneqq\zeta^{4i}\,C_0/L\in\bbQ(\zeta)\fps q,
\]
so that $(\Delta_i^{1/2})^2=\zeta^{8i}C_0^2/L^2=\Delta_i$. At $q=0$, $\Delta_i(0)=W^0_i\coloneqq\prod_{j\neq i}(\lambda_i-\lambda_j)/(5\lambda_i)=\zeta^{3i}$. The \emph{topological field theory of the formal quintic} is the TFT $\TFT$ on $\bbH\otimes\bbQ(\zeta)\fps q$ with
\[
\TFT_{g,n}(\tilde e_i,\dots,\tilde e_i)\coloneqq\Delta_i^{(2g-2+n)/2}=(\zeta^{4i}C_0/L)^{2g-2+n},
\]
and $\TFT_{g,n}$ vanishing on mixed colours; its unit is $1_{\TFT}\coloneqq\sum_i\Delta_i^{-1/2}\tilde e_i=\sum_ie_i=1$. We also regard $\TFT$ as a TFT on the colour space, with the orthonormal basis $\tilde e_i$ ($i\in\bbZ/5$) and the unit $1_{\TFT}$. By \cref{thm:set-GT}, since $R(0)=1$, $\TFT$ is the degree-zero part of the cohomological field theory along the line.
\end{defn}

% Section 2, part 4 of 8: the R-matrix and its rows; (QRR); the twist sign; the derivation of (QRR); h_w rotation; the graph sum; the two forms of the translation; (T).

%-------------------------------------------------------------------------------
\subsection{The \texorpdfstring{$R$}{R}-matrix, its rows and quantum Riemann--Roch}
\label{sec:set-rows}

\begin{defn}[The $R$-matrix and its rows]
\label{def:set-rows}
Let $R(z)=1+\sum_{k\geq1}R_kz^k$, with $R_k\in\End(\bbH)\otimes\bbQ(\zeta)\fps q$, be the $R$-matrix of the formal quintic along the line $t=T(q)H$: the $R$-matrix of its semisimple cohomological field theory \cite[Section~4.1]{GJR18}, in Givental's convention. This means:
\begin{itemize}
\item[(a)] $R$ is symplectic: $R(-z)^T\eta\,R(z)=\eta$, as matrices in the flat basis;
\item[(b)] for every colour $i$ the column $R(z)\tilde e_i\re^{u_i/z}$ solves the quantum differential equation $z\partial_Tv=H*v$, that is, $w\coloneqq R(z)\tilde e_i$ satisfies $z\partial_Tw+\mu_iw=H*w$ (\cref{cor:set-cyclicPF5}~(a)); thus the fundamental solution has the form $R(z)\re^{u/z}$ in the normalised canonical basis;
\item[(c)] in the graph sums of \cref{def:set-graphsum}, $R$ acts with $R^{-1}(\psi)$ on the legs.
\end{itemize}
Its value at $q=0$ is the quoted input (QRR), \cref{thm:set-qrr}, and its graph sum is the cohomological field theory of the formal quintic, the quoted input (T), \cref{thm:set-GT}.

The \emph{rows} are the lowered components $\Qr^{(i)}_a(z)\coloneqq\eta(H^a,R(z)\tilde e_i)\in\bbQ(\zeta)\fps q\fps z$, with the coefficients $\Qr^{(i)}_{a,n}\coloneqq\coef{z^n}\Qr^{(i)}_a$, and $\Qr$ denotes the matrix $(\Qr^{(i)}_a)$, with rows indexed by $a$ and columns by $i$. With $\zh=L^4z$ we \emph{define}
\[
\Phi_i(\zh,u)\coloneqq\frac{C_0\,\Qr^{(i)}_0(z)}{\zeta^iL},\qquad\text{i.e.}\quad\coef{\zh^n}\Phi_i\coloneqq\frac{C_0\,\Qr^{(i)}_{0,n}}{\zeta^iL^{4n+1}},
\]
and the coefficients $\mathrm{ph}_n(u)$ by $\Phi_0(\zh,u)=\sum_{n\geq0}\mathrm{ph}_n(u)\,\zh^n$, power series in $q=(1-u)/5^5$; $\mathrm{ph}_0=1$, because $R(0)=1$ and $\eta(1,\tilde e_0)=\Delta_0^{-1/2}=L/C_0$ (\cref{def:set-frob}). The \emph{flat amplitudes} are
\[
A_n\coloneqq\coef{z^n}\bigl(C_0\,\Qr^{(0)}_0\bigr)=L^{4n+1}\,\mathrm{ph}_n(L^{-5}),\qquad\text{so that}\quad\sum_nA_nz^n=L\,\Phi_0(L^4z,L^{-5}).
\]
\end{defn}

\begin{proof}[Proof of \cref{prop:set-cyclicLR}~(iii) and of the end of \cref{cor:set-cyclicPF5}~(b)]
For each colour $i$ put $w\coloneqq R(z)\tilde e_i$; by \cref{def:set-rows}~(b) it satisfies $z\partial_Tw+\mu_iw=H*w$. Taking lowered components as in the proof of \cref{prop:set-cyclicLR}~(ii), $z\partial_Tw_a+\mu_iw_a=\eta(H^a,H*w)=\eta(H*H^a,w)=\kappa^{\mathrm{LR}}_aw_{a+1}$ with $w_a=\Qr^{(i)}_a$; multiplying by $C_1$, with $\Dq=C_1\partial_T$ and $C_1\mu_i=\zeta^iL$ (\cref{cor:set-cyclicPF5}~(a)), gives $(z\Dq+\zeta^iL)\Qr^{(i)}_a=C_1\kappa^{\mathrm{LR}}_a\Qr^{(i)}_{a+1}$. Dividing by $\zeta^i$ gives the relations for $X_a$, with $C_1\kappa^{\mathrm{LR}}_4=C_0$ in the closing relation, and $X_0=\zeta^{-i}\Qr^{(i)}_0=L\Phi_i/C_0$ is the definition of $\Phi_i$. For the WKB form, the column $v=R(z)\tilde e_i\re^{u_i/z}$, with $\re^{u_i/z}$ a formal exponential, $\Dq\re^{u_i/z}=z^{-1}(\Dq u_i)\re^{u_i/z}$, solves $z\partial_Tv=H*v$ by \cref{def:set-rows}~(b); so \cref{cor:set-cyclicPF5}~(b) applies to it, and $C_0v_0=\re^{u_i/z}C_0\Qr^{(i)}_0=\zeta^i\re^{u_i/z}L\Phi_i$. Multiplying by the constant $z\zeta^{-i}$ gives the statement.
\end{proof}

The colour covariance $\Qr^{(i)}_{a,n}=\zeta^{i(a+1-n)}\Qr^{(0)}_{a,n}$, equivalently $\Phi_i(\zh,u)=\Phi_0(\zeta^{-i}\zh,u)$, is \cref{lem:gau-colrows}; the equation satisfied by $\Phi_0$ and its free coefficients are \cref{lem:gau-HPF}.

\begin{thm}[The $R$-matrix at $q=0$, (QRR); quoted]
\label{thm:set-qrr}
For a finite multiset $\mathbf w$ of non-zero numbers put
\begin{equation}
\label{eq:set-hw}
\hw(z)\coloneqq\exp\Bigl(-\sum_{k\geq1}\bk_k\,z^{2k-1}\sum_{a\in\mathbf w}a^{1-2k}\Bigr),\qquad\bk_k\coloneqq\frac{B_{2k}}{2k(2k-1)},
\end{equation}
and $\bw_i\coloneqq\{\zeta^i-\zeta^j\}_{j\neq i}\cup\{-5\zeta^i\}$. In Givental's convention (\cref{def:set-rows}), $R(q=0)=\diag(h_{\bw_0},\dots,h_{\bw_4})$ in the basis $\tilde e_i(q=0)$: $R(0,z)\tilde e_i(0)=h_{\bw_i}(z)\,\tilde e_i(0)$, where $R(0,z)$ is the part of order $q^0$ of $R(z)$, to all orders in $z$.
\end{thm}

This is \cite[proof of Lemma~5.2, (17)]{GJR18} and \cite[proof of Lemma~4.6, (19)--(20)]{GJR17} at $t=0$; both refer to \cite[Section~2.2, Theorem~1 and the Remark after it]{Giv01} and to \cite[Theorem~1, Corollary~1]{MR2276766}. It is the quoted input (QRR) of the note. It is used for every colour, and only in the following form.

\emph{Form used.} For every colour $i$,
\begin{equation}
\label{eq:set-qrr}
\Phi_i(z,1)=h_{\bw_i}(z);
\end{equation}
in particular, with $h\coloneqq h_{\bw_0}$ and $h_n\coloneqq\coef{z^n}h$, $A_n(q=0)=\mathrm{ph}_n(1)=h_n$. (In \cref{sec:gau} the colour-$0$ column of the actual $R$-matrix is written $\Phiact$, to distinguish it from the canonical one.)

\begin{proof}[Proof of the form used]
The point $q=0$ is $u=1$, where $L=1$, $\zh=z$ and $C_0=1$. For every $q$, $\eta(1,\tilde e_i)=\Delta_i^{1/2}\eta(1,e_i)=\Delta_i^{1/2}\eta(e_i,e_i)=\Delta_i^{-1/2}=\zeta^{-4i}L/C_0=\zeta^iL/C_0$ (\cref{cor:set-cyclicPF5}~(a) and the square roots of \cref{def:set-frob}). By (QRR), $R(0,z)\tilde e_i(0)=h_{\bw_i}(z)\tilde e_i(0)$, so the part of order $q^0$ of $\Qr^{(i)}_0$ is $h_{\bw_i}(z)\,\eta(1,\tilde e_i)|_{q=0}=\zeta^ih_{\bw_i}(z)$, and $\Phi_i(z,1)=\zeta^{-i}\Qr^{(i)}_0|_{q=0}=h_{\bw_i}(z)$.
\end{proof}

\begin{rmk}[The twist sign]
\label{rmk:set-qrrsign}
The weight $-5\zeta^i$ in $\bw_i$ is minus the weight $5\zeta^i$ of the fibre of $\cO(5)$ at $p_i$: this is the sign of a twist by an Euler class, for which the Chern root $5\zeta^i$ contributes the factor $\exp\bigl(+\sum_m\bk_m(z/5\zeta^i)^{2m-1}\bigr)=\exp\bigl(-\sum_m\bk_mz^{2m-1}(-5\zeta^i)^{1-2m}\bigr)$. The same sign appears in \cite[(17)]{GJR18}, in \cite[Corollary~1]{MR2276766} and in the vertex class \cite[(26)]{LP20} (\cref{rmk:set-qrrderiv}, Step~2). Pairing $j$ with $5-j$ gives $\frac1{1-\zeta^j}+\frac1{1-\zeta^{-j}}=1$, so $\sum_{j=1}^4\frac1{1-\zeta^j}=2$ and $\sum_{a\in\bw_0}a^{-1}=2-\frac15=\frac95$; hence $h_1=-\bk_1\cdot\frac95=-\frac1{12}\cdot\frac95=-\frac3{20}$, whereas the weight $+5$ would give $-\frac{11}{60}$. That $\mathrm{ph}_1=-\frac3{20}$ is forced, independently of (QRR), is part of \cref{lem:gau-HPF}.
\end{rmk}

\begin{rmk}[How (QRR) follows from the cohomological field theory]
\label{rmk:set-qrrderiv}
This remark is not used by the proofs of the note. It derives $R(q=0)$ from the cohomological field theory, assuming (T) and two classical facts on the Hodge bundle $\bbE$. The letter $\mathsf h$ is local.

\emph{Step 1 (the quantum differential equation alone).} $R(0,z)$ is diagonal in the basis $\tilde e_i(0)$, with entries $c_i(z)\in\bbQ(\zeta)\fps z$, $c_i(0)=1$, $c_i(z)c_i(-z)=1$. Indeed, the $R$-matrix of \cite[Theorem~4.1]{GJR18} has the form $\mathbf E\exp h$ of \cite[8.7~Proposition~(ii)]{MR2917177}, with $\mathbf E$ a solution of the system of \cite[8.2~Proposition~(ii)]{MR2917177} and $\exp h$ constant and diagonal, so its columns solve the quantum differential equation, as in \cref{def:set-rows}~(b). In the form $(z\Dq+\Dq u_i)(R\tilde e_i)=C_1\,H*(R\tilde e_i)$ take the part of order $q^0$: $\Dq$ kills it, $\Dq u_i(0)=\zeta^i$, $C_1(0)=1$, and $H*$ at $q=0$ is the classical product. So $R(0,z)\tilde e_i(0)$ lies, coefficientwise in $z$, in the $\zeta^i$-eigenline of the classical multiplication by $H$, which is spanned by $\tilde e_i(0)$, the eigenvalues $\zeta^0,\dots,\zeta^4$ being distinct. The symplectic condition \cite[(5)]{PPZ15}, with $\eta(\tilde e_i,\tilde e_j)=\delta_{ij}$, gives $c_i(z)c_i(-z)=1$; so $c_i$ is the exponential of an odd series.

\emph{Step 2 (the theory at $q=0$).} The part of order $q^0$ comes from the constant maps: $\overline{\mathcal M}_{g,n}(\bbP^4,0)=\overline{\mathcal M}_{g,n}\times\bbP^4$, the virtual class is $e(\bbE^\vee\boxtimes T_{\bbP^4})$ and $R\pi_*f^*\cO(5)=\cO(5)-\bbE^\vee\otimes\cO(5)$. In degree $0$ the insertions of $TH$ give classes pulled back along forgetful maps, whose push-forwards vanish, and the divisor equation holds for the classes; so the part of order $q^0$ at $t=TH$ is the degree-$0$ part. Localisation on $\bbP^4$ gives $\Omega_{g,n}(\gamma_1,\dots,\gamma_n)|_{q=0}=\sum_j\prod_k\gamma_k(p_j)\,H_g^{(p_j)}$ with \cite[(26)]{LP20}
\[
H_g^{(p_i)}=\frac{e(\bbE^*\otimes T_{p_i}\bbP^4)}{e(T_{p_i}\bbP^4)}\cdot\frac{5\lambda_i}{e(\bbE^*\otimes(5\lambda_i))}.
\]
Write $\mathsf h(w)\coloneqq e(\bbE^*\otimes\bbC_w)$ for the line $\bbC_w$ of weight $w$. Mumford's relation in the form $\mathsf h(w)\mathsf h(-w)=(-1)^gw^{2g}$ \cite[Section~2.1]{MR1728879} gives $H_g^{(p_i)}=W_i^{g-1}\prod_{a\in\bw_i}\mathsf h(a)/a^g$ with $W_i=\prod_{j\neq i}(\lambda_i-\lambda_j)/(5\lambda_i)=\Delta_i(0)$, and $\mathsf h(a)/a^g=\exp\bigl(-\sum_k(2k-2)!\,\mathrm{ch}_{2k-1}(\bbE)\,a^{1-2k}\bigr)$, using $\mathrm{ch}_{2k}(\bbE)=0$ for $k\geq1$. So in the basis $\tilde e_i(0)$, $\Omega|_{q=0}=\bigoplus_i\Omega^{(i)}$, where $\Omega^{(i)}_{g,n}(\tilde e_i,\dots,\tilde e_i)=\Delta_i(0)^{(2g-2+n)/2}\exp\bigl(\sum_ks^{(i)}_k\mathrm{ch}_{2k-1}(\bbE)\bigr)$, $s^{(i)}_k\coloneqq-(2k-2)!\sum_{a\in\bw_i}a^{1-2k}$, and mixed colours give $0$.

\emph{Step 3 (Hodge classes as $R$-matrix actions).} For an odd series $h(z)=\sum_jh_{2j-1}z^{2j-1}$, the theory with the classes $\exp\{\sum_jh_{2j-1}\frac{(2j)!}{B_{2j}}\mathrm{ch}_{2j-1}(\bbE)\}$ times a rank-one TFT $\omega$ \cite[(8.6)]{MR2917177} is the transform $T_z^{-1}\circ\exp h(z)\circ T_z$ of $\omega$ \cite[8.7~Proposition~(i)]{MR2917177}. In the convention of \cite[Def.~2.13]{PPZ15} this is $\exp(h).\omega$: \cite[p.~20]{PPZ15} use the inverse $R$-matrix on the legs and state that Givental and Teleman use the same convention, and the sign is confirmed in genus one: for $R=\exp(\sum_kr_{2k-1}z^{2k-1})$ acting on the trivial rank-one theory, $\int_{\overline{\mathcal M}_{1,1}}(R.\omega)_{1,1}(1)=r_1(\int\kappa_1-\int\psi_1)+\frac{r_1}2=\frac{r_1}2$, the last term coming from the graph with one loop ($|\Aut|=2$, edge term $r_1$), while $\int_{\overline{\mathcal M}_{1,1}}a_1\lambda_1=a_1/24$ for the class $\exp(a_1\mathrm{ch}_1(\bbE)+\dots)$; so $r_1=a_1/12=a_1B_2/2!$, as $h_{2k-1}=a_kB_{2k}/(2k)!$ requires, whereas opposite conventions would change the sign of every odd coefficient. Matching $h_{2k-1}(2k)!/B_{2k}=s^{(i)}_k$ gives $h_{2k-1}=-\bk_k\sum_{a\in\bw_i}a^{1-2k}$, i.e.\ $\exp h=h_{\bw_i}$ by \eqref{eq:set-hw}; so $\Omega^{(i)}=h_{\bw_i}.\omega^{(i)}$, $\omega^{(i)}$ the rank-one TFT of colour $i$. For colour $0$, $a_1=-\sum_{a\in\bw_0}a^{-1}=-\frac95$ and $r_1=-\frac3{20}=h_1$ (\cref{rmk:set-qrrsign}).

\emph{Step 4 (the graph sum at $q=0$).} By (T), $\Omega=R.\TFT$ (\cref{thm:set-GT}). For fixed $(g,n)$ and cohomological degree the graph sum is finite and polynomial in the coefficients of $R$ and $\TFT$, so taking the part of order $q^0$ commutes with it, and $\Omega|_{q=0}=R(0).\TFT(0)=\bigoplus_ic_i.\omega^{(i)}$ by Step~1: a diagonal $R$-matrix acting on a direct sum of rank-one TFTs keeps the colours, since a vertex term vanishes unless its half-edges have one colour and the edges and the translation are diagonal.

\emph{Step 5 (uniqueness in rank one).} By Steps 2--4, $c_i.\omega^{(i)}=h_{\bw_i}.\omega^{(i)}$. With $\delta_i\coloneqq c_i/h_{\bw_i}=\exp(\sum_kd_kz^{2k-1})$ and the group property \cite[Prop.~2.14]{PPZ15}, $\delta_i.\omega^{(i)}=\omega^{(i)}$, i.e.\ by Step~3 $\exp\bigl(\sum_kd_k\frac{(2k)!}{B_{2k}}\mathrm{ch}_{2k-1}(\bbE)\bigr)=1$ on every $\overline{\mathcal M}_{g,n}$. Taking logarithms (the class is unipotent) and separating degrees, $d_k\frac{(2k)!}{B_{2k}}\mathrm{ch}_{2k-1}(\bbE)=0$. If for every $k\geq1$ the class $\mathrm{ch}_{2k-1}(\bbE)$ is non-zero in $H^*(\overline{\mathcal M}_{g,n})$ for some $(g,n)$, then $d_k=0$ for all $k$, $c_i=h_{\bw_i}$, and $R(q=0)=\diag(h_{\bw_i})$, which is (QRR). The rank-one theories are described in \cite[8.10~Proposition]{MR2917177}; that description is not used here.\draftnote{Mumford's relation $\mathrm{ch}_{2k}(\bbE)=0$ and, for every $k\geq1$, the non-vanishing of $\mathrm{ch}_{2k-1}(\bbE)$ in some $H^*(\overline{\mathcal M}_{g,n})$ are used here without a pinned statement}
\end{rmk}

\begin{lem}[Rotation of the colours]
\label{lem:set-hrot}
$h_{\bw_i}(z)=h_{\bw_0}(\zeta^{-i}z)$ for every $i$.
\end{lem}

\begin{proof}
As multisets $\bw_i=\zeta^i\bw_0$, since $\zeta^i-\zeta^j=\zeta^i(1-\zeta^{j-i})$ and $-5\zeta^i=\zeta^i(-5)$. Hence $\sum_{a\in\bw_i}a^{1-2k}=\zeta^{i(1-2k)}\sum_{a\in\bw_0}a^{1-2k}$, and $z^{2k-1}\zeta^{i(1-2k)}=(\zeta^{-i}z)^{2k-1}$ in \eqref{eq:set-hw}.
\end{proof}

%-------------------------------------------------------------------------------
\subsection{The graph sum}
\label{sec:set-graph}

\begin{defn}[The Givental--Teleman graph sum]
\label{def:set-graphsum}
Let $\Omega$ be a cohomological field theory with unit $1$ \cite[Def.~0.3]{PPZ15} on a free module $V$ of finite rank with pairing $\eta$ over a commutative $\bbQ$-algebra, and let $R(z)=1+O(z)\in\End(V)\fps z$ be symplectic, $R(-z)^*R(z)=1$ with ${}^*$ the adjoint for $\eta$ \cite[(5)]{PPZ15}. The ancestor correlators are $\langle\psi^{k_1}a_1\cdots\psi^{k_n}a_n\rangle_h\coloneqq\int_{\overline{\mathcal M}_{h,n}}\Omega_{h,n}(a_1,\dots,a_n)\prod_j\psi_j^{k_j}$; for $q\in V\fps z$ (a point of the loop space, in Givental's notation; not the parameter $q$) the dilaton-shifted potential is $\overline F^\Omega_h(q)\coloneqq\sum_n\frac1{n!}\langle(q+z1)^n\rangle_h$, a point being inserted through $z^ka\mapsto\psi^ka$ (base point $q=-z1$), and $\dd^n\overline F^\Omega_h(q)$ is its $n$-th derivative, the symmetric $n$-linear form $(x_1,\dots,x_n)\mapsto\sum_m\frac1{m!}\langle x_1\cdots x_n(q+z1)^m\rangle_h$ on $V\fps z$. The \emph{edge} and the \emph{leaf point} of $R$ are
\begin{gather*}
E_R(\psi,\psi')\coloneqq\frac{\eta^{-1}-\bigl(R^{-1}(\psi)\otimes R^{-1}(\psi')\bigr)\eta^{-1}}{\psi+\psi'}\in(V\otimes V)\fps{\psi,\psi'},\\
q_R(z)\coloneqq R^{-1}(z)(-z1)=-z1+\ttr_R(z),
\end{gather*}
where $\eta^{-1}\in V\otimes V$ is the inverse bivector, $E_R$ is a power series by the symplectic condition \cite[Rem.~2.1, (6)]{PPZ15}, and $\ttr_R(z)\coloneqq z1-zR^{-1}(z)1\in z^2V\fps z$ is the translation of \cite[Def.~2.13]{PPZ15}. For a stable graph $\Gamma$ without legs, with vertices $v$ of genus $g_v$ and valence $n_v$ and of total genus $g=\sum_vg_v+h^1(\Gamma)$, the contraction
\[
\mathrm{Cont}_\Gamma[E,q;\Omega]\coloneqq|\Aut\Gamma|^{-1}\Bigl\langle\prod_v\dd^{n_v}\overline F^\Omega_{g_v}(q),\ \prod_eE\Bigr\rangle
\]
contracts, at each vertex, the $n_v$-th derivative of $\overline F^\Omega_{g_v}$ at $q$, one slot for each half-edge, with one copy of the bivector $E(\psi,\psi')$ for each edge, its two slots at the two half-edges of that edge; it includes the factor $|\Aut\Gamma|^{-1}$, so that sums over graphs carry no further factor. For $q=q_R$ the contraction is a finite sum, as $\ttr_R\in z^2V\fps z$. Let $R.\Omega$ be the unit-preserving action of \cite[Defs.~2.2, 2.5, 2.13]{PPZ15}: legs $R^{-1}(\psi)$, edges $E_R$, translation $\ttr_R$; \cite[p.~20]{PPZ15} use the inverse $R^{-1}$ on the legs, the convention of Givental and Teleman. The translation of \cite[Def.~2.5]{PPZ15} pushes the translation points forward along the forgetful maps, so that the half-edges carry the classes $\psi$ of $\overline{\mathcal M}_{g_v,n_v}$, whereas the vertex tensors above insert these points and use the classes $\psi$ of $\overline{\mathcal M}_{g_v,n_v+m}$. The two forms agree: for $g\geq2$, $\mathrm{Cont}_\Gamma[E_R,q_R;\Omega]$ is the term of $\Gamma$ in $\int_{\overline{\mathcal M}_g}(R.\Omega)_{g,0}$ (\cref{lem:set-pushins}~(iii) below). The sum $\sum_\Gamma\mathrm{Cont}_\Gamma[E_R,q_R;\Omega]$ over the leg-free stable graphs of genus $g$ is the \emph{graph sum of genus $g$}. The definitions and propositions of \cite{PPZ15} are stated over $\bbQ$; those used in the note are identities between finite sums, polynomial in the coefficients of the $R$-matrices and of the cohomological field theory, and they hold verbatim over every commutative $\bbQ$-algebra.

\emph{Colours.} When $\Omega=\TFT$ (\cref{def:set-frob}), which vanishes on mixed colours, a \emph{coloured graph} is a stable graph with a colour at each vertex, and the term of $\Gamma$ is the sum over its colourings. For a constant isometry $A\colon(W,\eta_W)\to(V,\eta)$ and $A_*\Omega\coloneqq\Omega\circ(A^{-1})^{\otimes n}$ one has $A_*(R.\Omega)=(ARA^{-1}).(A_*\Omega)$, which is immediate from \cite[Defs.~2.2, 2.13]{PPZ15}; it gives the graph sum in the colour basis, whose edges and leaves are written out in \cref{sec:sgn}.
\end{defn}

\begin{lem}[The translation: push-forward and insertion]
\label{lem:set-pushins}
Let $\Omega$ be a cohomological field theory \cite[Def.~0.3]{PPZ15} on a free module $V$ of finite rank with pairing $\eta$ over a commutative $\bbQ$-algebra, with the correlators $\langle\cdots\rangle_h$ of \cref{def:set-graphsum}. Let $2h-2+n>0$, $m\geq0$, and let $p_m\colon\overline{\mathcal M}_{h,n+m}\to\overline{\mathcal M}_{h,n}$ forget the last $m$ points.
\begin{enumerate}
\item For $k_1,\dots,k_n\geq0$ and $c_1,\dots,c_m\geq1$, in $H^*(\overline{\mathcal M}_{h,n+m})$,
\[
\prod_{j=1}^n\psi_j^{k_j}\prod_{s=1}^m\psi_{n+s}^{c_s}=\prod_{j=1}^n(p_m^*\psi_j)^{k_j}\prod_{s=1}^m\psi_{n+s}^{c_s}.
\]
\item Let $\ttr\in z^2V\fps z$ and let $\ttr\Omega$ be the translation of $\Omega$ by $\ttr$ \cite[Def.~2.5]{PPZ15}. For $a_1,\dots,a_n\in V$ and $k_1,\dots,k_n\geq0$,
\[
(\ttr\Omega)_{h,n}(a_1,\dots,a_n)\prod_{j=1}^n\psi_j^{k_j}=\sum_{m\geq0}\frac1{m!}\,p_{m*}\Bigl(\Omega_{h,n+m}\bigl(a_1,\dots,a_n,\ttr(\psi_{n+1}),\dots,\ttr(\psi_{n+m})\bigr)\prod_{j=1}^n\psi_j^{k_j}\Bigr),
\]
where on the right $\psi_j$ is the class of $\overline{\mathcal M}_{h,n+m}$; consequently
\[
\int_{\overline{\mathcal M}_{h,n}}(\ttr\Omega)_{h,n}(a_1,\dots,a_n)\prod_{j=1}^n\psi_j^{k_j}=\sum_{m\geq0}\frac1{m!}\bigl\langle\psi^{k_1}a_1\cdots\psi^{k_n}a_n\,\ttr^m\bigr\rangle_h ,
\]
$\ttr$ being inserted at $m$ further points through $z^ca\mapsto\psi^ca$. The sums are finite. If $\Omega$ has the unit $1$, the right side is $\dd^n\overline F^\Omega_h(q)(\psi^{k_1}a_1,\dots,\psi^{k_n}a_n)$ at $q=-z1+\ttr$.
\item Let $\Omega$ have the unit $1$ and let $R(z)=1+O(z)$ be symplectic, as in \cref{def:set-graphsum}. For $g\geq2$ and every leg-free stable graph $\Gamma$ of genus $g$, $\mathrm{Cont}_\Gamma[E_R,q_R;\Omega]$ is the term of $\Gamma$ in $\int_{\overline{\mathcal M}_g}(R.\Omega)_{g,0}$, that is, $|\Aut\Gamma|^{-1}$ times the integral over $\overline{\mathcal M}_g$ of the contribution of $\Gamma$ in \cite[Def.~2.2]{PPZ15} to the action of $R$ on $\ttr_R\Omega$, which is $R.\Omega$ \cite[Def.~2.13]{PPZ15}. Consequently $\int_{\overline{\mathcal M}_g}(R.\Omega)_{g,0}$ is the graph sum of genus $g$. By \cite[Prop.~2.12]{PPZ15}, $R.\Omega$ is again a cohomological field theory with the unit $1$, to which (iii) applies.
\end{enumerate}
\end{lem}

\begin{proof}
(i) Let first $m=1$. For $N\geq0$ with $2h-2+N>0$ let $\pi\colon\overline{\mathcal M}_{h,N+1}\to\overline{\mathcal M}_{h,N}$ forget the last point, and for $j\leq N$ let $\delta_j$ be the class of the divisor of $\overline{\mathcal M}_{h,N+1}$ of the curves with a genus-$0$ component carrying only the markings $j$ and $N+1$ ($N$, $\pi$ and $\delta_j$ are local letters). The formula \cite[(8)]{PPZ15} for the map forgetting one point, $\pi^*(\psi_j^d)=\psi_j^d-\delta_j\,\pi^*(\psi_j^{d-1})$, gives $\pi^*\psi_j=\psi_j-\delta_j$ for $d=1$; for $d=2$ its right side is then $\psi_j^2-\delta_j\psi_j+\delta_j^2$, while $\pi^*(\psi_j^2)=(\psi_j-\delta_j)^2$, so $\psi_j\delta_j=0$. The automorphism of $\overline{\mathcal M}_{h,N+1}$ exchanging the markings $j$ and $N+1$ preserves the divisor and exchanges $\psi_j$ and $\psi_{N+1}$ under pull-back; hence $\psi_{N+1}\delta_j=0$ as well, that is, $\psi_j\psi_{N+1}=\pi^*(\psi_j)\,\psi_{N+1}$ for every $j\leq N$. It follows that, for $k_1,\dots,k_N\geq0$ and $c\geq1$,
\[
\textstyle\prod_{i=1}^N\psi_i^{k_i}\,\psi_{N+1}^c=\prod_{i=1}^N(\pi^*\psi_i)^{k_i}\,\psi_{N+1}^c,\tag{$*$}
\]
the classes $\psi_i$ on the left being those of $\overline{\mathcal M}_{h,N+1}$: replace the factors $\psi_i$ by their pull-backs one at a time, by the identity just proved, which needs only one factor $\psi_{N+1}$. For $N=n$, $\pi=p_1$ and $(*)$ is (i) for $m=1$. Let now $m\geq2$ and assume (i) for $m-1$. Write $p_m=p_{m-1}\circ\pi$, where $\pi\colon\overline{\mathcal M}_{h,n+m}\to\overline{\mathcal M}_{h,n+m-1}$ forgets the last point, and let $Y\coloneqq\prod_{s<m}\psi_{n+s}^{c_s}$ and $X\coloneqq\prod_j\psi_j^{k_j}\,Y$ (local letters) be these monomials in the classes of $\overline{\mathcal M}_{h,n+m-1}$. By $(*)$ for $\pi$, with $N=n+m-1$ and $c=c_m$, the left side of (i) is $\pi^*(X)\,\psi_{n+m}^{c_m}$. By (i) for $m-1$, $X=\prod_j(p_{m-1}^*\psi_j)^{k_j}\,Y$, so $\pi^*(X)=\prod_j(p_m^*\psi_j)^{k_j}\,\pi^*(Y)$; and by $(*)$ again, $\pi^*(Y)\,\psi_{n+m}^{c_m}=\prod_{s<m}\psi_{n+s}^{c_s}\,\psi_{n+m}^{c_m}$ in the classes of $\overline{\mathcal M}_{h,n+m}$. Together, the left side of (i) is the right side.

(ii) By \cite[Def.~2.5]{PPZ15} and the convention stated after it, $(\ttr\Omega)_{h,n}(a_1,\dots,a_n)=\sum_m\frac1{m!}\,p_{m*}Z_m$ with $Z_m\coloneqq\Omega_{h,n+m}(a_1,\dots,a_n,\ttr(\psi_{n+1}),\dots,\ttr(\psi_{n+m}))$ (a local letter), where an argument $\ttr(\psi_{n+s})$ stands for $\sum_{c\geq2}\psi_{n+s}^c$ times the class with the argument $\ttr_c\coloneqq\coef{z^c}\ttr$. By the projection formula, $p_{m*}(Z_m)\prod_j\psi_j^{k_j}=p_{m*}\bigl(Z_m\prod_j(p_m^*\psi_j)^{k_j}\bigr)$. Each term of $Z_m$ is a class $\Omega_{h,n+m}(a_1,\dots,a_n,\ttr_{c_1},\dots,\ttr_{c_m})$ times $\prod_s\psi_{n+s}^{c_s}$ with every $c_s\geq2$, so by (i) the pulled-back classes $p_m^*\psi_j$ may be replaced by the classes $\psi_j$ of $\overline{\mathcal M}_{h,n+m}$; this is the first identity. The second follows by integration, since $\int_{\overline{\mathcal M}_{h,n}}p_{m*}(W)=\int_{\overline{\mathcal M}_{h,n+m}}W$ for every class $W$: the term $m$ becomes $\frac1{m!}\langle\psi^{k_1}a_1\cdots\psi^{k_n}a_n\,\ttr^m\rangle_h$. A term with $m$ translation points has degree at least $2m$ on $\overline{\mathcal M}_{h,n+m}$, of dimension $3h-3+n+m$, so it vanishes for $m>3h-3+n$, and only the $c_s\leq3h-3+n+m$ contribute. The last sentence is the definition of $\dd^n\overline F^\Omega_h$ in \cref{def:set-graphsum}, as $q+z1=\ttr$.

(iii) By \cite[Def.~2.13]{PPZ15}, $R.\Omega$ is the action of \cite[Def.~2.2]{PPZ15} on the translation of $\Omega$ by $\ttr_R(z)=z1-zR^{-1}(z)1\in z^2V\fps z$. In \cite[Def.~2.2]{PPZ15}, $(R.\Omega)_{g,0}=\sum_\Gamma|\Aut\Gamma|^{-1}C_\Gamma$ (a local letter) over the leg-free stable graphs of genus $g$, where $C_\Gamma$ is the push-forward, along the gluing map $\prod_v\overline{\mathcal M}_{g_v,n_v}\to\overline{\mathcal M}_g$ of \cite[\S0.2, (1)]{PPZ15}, of the class obtained by placing $(\ttr_R\Omega)_{g_v,n_v}$ at each vertex and, at each edge, the bivector $E_R(\psi',\psi'')$ in the classes $\psi'$, $\psi''$ of its two half-edges on the moduli spaces of their vertices, and contracting; $E_R$ is the edge of \cite[Def.~2.2~(iii)]{PPZ15}, which, as all formulas of that definition, is formed with $R^{-1}$ \cite[p.~20]{PPZ15}. Since the push-forward preserves integrals, the term of $\Gamma$ is $|\Aut\Gamma|^{-1}$ times the integral of that class over $\prod_v\overline{\mathcal M}_{g_v,n_v}$. Expand every edge class into monomials $\psi'^{\,a}\psi''^{\,b}$ with coefficients in $V\otimes V$; only finitely many of them have degree at most the dimension. For each choice of monomials the integral over the product is the contraction, with these coefficients, of the tensor product over the vertices of the multilinear forms $(x_1,\dots,x_{n_v})\mapsto\int_{\overline{\mathcal M}_{g_v,n_v}}(\ttr_R\Omega)_{g_v,n_v}(x_1,\dots,x_{n_v})\prod_l\psi_l^{k_l}$, where $k_l$ is the exponent at the half-edge $l$ of $v$. By (ii), as $q_R+z1=\ttr_R$, this form is $\dd^{n_v}\overline F^\Omega_{g_v}(q_R)(\psi^{k_1}x_1,\dots,\psi^{k_{n_v}}x_{n_v})$. Summing over the monomials gives $|\Aut\Gamma|^{-1}\langle\prod_v\dd^{n_v}\overline F^\Omega_{g_v}(q_R),\prod_eE_R\rangle=\mathrm{Cont}_\Gamma[E_R,q_R;\Omega]$, and summing over $\Gamma$ gives the graph sum.
\end{proof}

\begin{thm}[The reconstruction (T); quoted]
\label{thm:set-GT}
The cohomological field theory of the formal quintic along the line, i.e.\ the classes $\Omega_{g,n}$ of \cref{def:set-fq} with $\Qn=\re^{T(q)}$ (which, by the divisor equation, is the theory at the point $t=T(q)H$ with $\Qn=1$), equals $R.\TFT$, with $R$ of \cref{def:set-rows} and $\TFT$ of \cref{def:set-frob} \cite[Theorem~(2) in (1.15)]{MR2917177}, \cite[Section~4.1, Theorem~4.1]{GJR18}. In particular, for $g\geq2$, since $F_g=\int_{\overline{\mathcal M}_g}\Omega_{g,0}$ (\cref{def:set-fq}), \cref{lem:set-pushins}~(iii) gives
\[
F_g=\sum_\Gamma\mathrm{Cont}_\Gamma[E_R,q_R;\TFT]
\]
as power series in $q$, the sum running over the leg-free stable graphs of genus $g$.
\end{thm}

This is the quoted input (T); its second assertion follows from the first by \cref{lem:set-pushins}~(iii), applied to $\TFT$ with its unit $1$ and to the symplectic series $R$ with $R(0)=1$ (\cref{def:set-rows}~(a)). In Theorems~A and~C it identifies the $R$-matrix whose rows enter the outer graph sum (\cref{sec:set-hat}) and define the normalisation series $\Dn$ (\cref{thm:gau-R0}); in Corollary~B it identifies the total with $F_g$ (\cref{cor:mai-gap}).

% Section 2, part 5 of 8: the Frobenius basis at the conifold point; conifold frames; the differential ring and the conifold expansion; tau; the canonical coordinates.

%-------------------------------------------------------------------------------
\subsection{Conifold frames and the conifold expansion}
\label{sec:set-frames}

In this subsection functions of $u$ are formal series in $u$ (or in $s=u^{1/5}$), with $\Dq=-(1-u)\,\dd/\dd u$ (\cref{def:set-points}), and $\PFfour$ is taken in its $u$-form (\cref{def:set-genus0}).

\begin{lem}[The Frobenius basis at the conifold point]
\label{lem:set-frobbasis}
Let $\theta\coloneqq u\,\dd/\dd u$, so that $\Dq=(1-u^{-1})\theta$ on $\bbQ\lser u$, and let $\ell$ be an indeterminate with $\Dq\ell=1-u^{-1}$, i.e.\ $\theta\ell=1$, a formal $\log u$ (the indeterminate of \cref{def:gau-conc}).
\begin{enumerate}
\item $\PFfour=\sum_{j=-3}^{1}u^jP_j(\theta)$ with polynomials $P_j$, where
\[
P_{-3}(x)=x(x-1)^2(x-2),\qquad P_{-2}(x)=-x(x-1)\bigl(4x^2-6x+\tfrac{17}5\bigr),\qquad P_{-1}(0)=0,\quad P_{-1}'(0)=-\tfrac75 .
\]
Hence $\PFfour$ maps $\bbQ\fps u$ into itself, its local exponents at $u=0$ are $0,1,1,2$, and a power-series solution is determined by its coefficients of $u^0,u^1,u^2$; $f_k$ ($k=0,1,2$) denotes the solution with $f_k=u^k+O(u^3)$.
\item For $g=\sum_nb_nu^n\in\bbQ\fps u$, $\PFfour(g\ell)=\ell\,\PFfour g+r(g)$ with $r(g)\coloneqq\sum_nb_n\sum_jP_j'(n)u^{n+j}\in u^{-3}\bbQ\fps u$. If $\PFfour g=0$, then $g\ell+h$ with $h\in\bbQ\fps u$ solves $\PFfour$ for some $h$ if and only if $\coef{u^m}r(g)=0$ for $m=-3,-2,-1$.
\item On $\operatorname{span}(f_0,f_1,f_2)$ the obstruction $(\coef{u^{-3}}r,\coef{u^{-2}}r,\coef{u^{-1}}r)$ has the matrix, in the columns $f_0,f_1,f_2$,
\[
\begin{pmatrix}-2&0&0\\ \frac{17}5&0&0\\ -\frac75&-\frac75&2\end{pmatrix},
\]
with kernel $\bbQ\,(0,1,\frac7{10})$. Hence, up to a non-zero factor, $\varpi_1\coloneqq f_1+\gamma f_2$, $\gamma\coloneqq\frac7{10}$, is the only power-series solution $g$ for which a solution $g\ell+h$, $h\in\bbQ\fps u$, exists; $h$ is unique up to adding power-series solutions, and the logarithmic solution with $\coef{u^0}h=\coef{u^1}h=\coef{u^2}h=0$ is $\varpi_1\ell+\rho_3$, $\rho_3\in u^3\bbQ\fps u$. The \emph{vanishing period} $\varpi_1=u+\frac7{10}u^2+O(u^3)$ is normalised by $\varpi_1=u+O(u^2)$.
\end{enumerate}
\end{lem}

For instance $f_0=1+\frac2{625}u^3+O(u^4)$, $f_1=u-\frac{19}{60}u^3+O(u^4)$, $f_2=u^2+\frac{37}{30}u^3+O(u^4)$ and $\varpi_1=u+\frac7{10}u^2+\frac{41}{75}u^3+O(u^4)$.

\begin{proof}
(i) Since $\theta u^m=u^m(\theta+m)$ and $\Dq=(1-u^{-1})\theta$, every operator $\sum_kc_k\Dq^k$ with $c_k\in\bbQ[u,u^{-1}]$ is uniquely of the form $\sum_ju^jP_j(\theta)$. Moving the factors $u^{-1}$ to the left with $\theta u^{-1}=u^{-1}(\theta-1)$, the part of order $u^{-k}$ of $\Dq^k$ is $(-1)^ku^{-k}\theta(\theta-1)\cdots(\theta-k+1)$, and its part of order $u^{1-k}$ is $(-1)^{k-1}u^{1-k}\theta(\theta-1)\cdots(\theta-k+2)\bigl(k\theta-\binom k2\bigr)$ (the word with the factor $\theta$ in position $p+1$ contributes $\theta(\theta-1)\cdots(\theta-k+2)\,(\theta-k+1+p)$, $p=0,\dots,k-1$). In the $u$-form $\PFfour=u\Dq^4-(1-u)\bigl(2\Dq^3+\frac75\Dq^2+\frac25\Dq+\frac{24}{625}\bigr)$, the terms of order $u^{-3}$ come from $u\Dq^4$ and $-2\Dq^3$, and those of order $u^{-2}$ from $u\Dq^4$, $-2(1-u)\Dq^3$ and $-\frac75\Dq^2$. This gives
\begin{align*}
P_{-3}(\theta)&=\theta(\theta-1)(\theta-2)(\theta-3)+2\theta(\theta-1)(\theta-2)=\theta(\theta-1)^2(\theta-2),\\
P_{-2}(\theta)&=-\theta(\theta-1)(\theta-2)(4\theta-6)-2\theta(\theta-1)(\theta-2)-6\theta(\theta-1)^2-\tfrac75\theta(\theta-1)\\
&=-\theta(\theta-1)\bigl(4\theta^2-6\theta+\tfrac{17}5\bigr).
\end{align*}
Every $\Dq^k$ with $k\geq1$ ends with the factor $\theta$, and the term $-\frac{24}{625}(1-u)$ contributes to $P_0$ and $P_1$ only; so $P_j(0)=0$ for $j\leq-1$. The value $P_{-1}'(0)$ is computed in (iii). Now $\PFfour u^n=\sum_jP_j(n)u^{n+j}$, and $P_j(n)=0$ whenever $n\geq0$ and $n+j<0$ (the roots $0,1,2$ of $P_{-3}$, $0,1$ of $P_{-2}$ and $0$ of $P_{-1}$); so $\PFfour$ maps $\bbQ\fps u$ into itself. For $g=\sum b_nu^n$ the coefficient of $u^N$, $N\geq0$, in $\PFfour g$ is $P_{-3}(N+3)b_{N+3}$ plus a combination of $b_{N-1},\dots,b_{N+2}$, and $P_{-3}(N+3)=(N+3)(N+2)^2(N+1)\neq0$. So the equation $\PFfour g=0$ determines $b_{N+3}$ from the lower coefficients, and $b_0,b_1,b_2$ are free; the indicial polynomial $P_{-3}$ has the roots $0,1,1,2$.

(ii) As $\theta\ell=1$, $\theta(u^n\ell)=u^n(n\ell+1)$, so $\theta$ acts on $\operatorname{span}(u^n\ell,u^n)$ as $n+N$ with $N(u^n\ell)=u^n$, $N(u^n)=0$, $N^2=0$. Hence $P(\theta)(u^n\ell)=u^n(P(n)\ell+P'(n))$ for every polynomial $P$, and $\PFfour(u^n\ell)=\ell\,\PFfour u^n+\sum_jP_j'(n)u^{n+j}$. Summing over $n$ (each power of $u$ receives finitely many contributions) gives $\PFfour(g\ell)=\ell\,\PFfour g+r(g)$. For $n\geq3$ all exponents $n+j$ are $\geq0$; so $r(g)\in u^{-3}\bbQ\fps u$, and its part of negative order depends only on $b_0,b_1,b_2$. If $\PFfour g=0$, then $\PFfour(g\ell+h)=r(g)+\PFfour h$ with $\PFfour h\in\bbQ\fps u$; a solution therefore requires $\coef{u^m}r(g)=0$ for $m<0$, and conversely, if these coefficients vanish, $\PFfour h=-r(g)$ is solved by the triangular recursion of (i), with $h_0,h_1,h_2$ free.

(iii) By (ii), $\coef{u^m}r(g)=\sum_jP_j'(m-j)\,b_{m-j}$, the sum over $j$ with $m-j\geq0$. For $m=-3$ this is $P_{-3}'(0)b_0=-2b_0$; for $m=-2$ it is $P_{-3}'(1)b_1+P_{-2}'(0)b_0=\frac{17}5b_0$, since $1$ is a double root of $P_{-3}$; for $m=-1$ it is $P_{-3}'(2)b_2+P_{-2}'(1)b_1+P_{-1}'(0)b_0=2b_2-\frac75b_1+P_{-1}'(0)b_0$. Finally $P_{-1}'(0)=\coef{u^{-1}}r(1)=\coef{u^{-1}}\PFfour(\ell)$. From $\Dq u^{-m}=mu^{-m-1}-mu^{-m}$ one gets $\Dq\ell=1-u^{-1}$, $\Dq^2\ell=u^{-1}-u^{-2}$, $\Dq^3\ell=-u^{-1}+3u^{-2}-2u^{-3}$ and $\Dq^4\ell=u^{-1}-7u^{-2}+12u^{-3}-6u^{-4}$, and the coefficient of $u^{-1}$ in $\PFfour\ell$ is $-7$ (from $u\Dq^4\ell$) $+\,(2+6)$ (from $-2(1-u)\Dq^3\ell$) $-\frac75-\frac75$ (from $-\frac75(1-u)\Dq^2\ell$) $+\frac25$ (from $-\frac25(1-u)\Dq\ell$), that is $-\frac75$. The power-series solutions are $b_0f_0+b_1f_1+b_2f_2$, which gives the matrix; its kernel consists of the $(b_0,b_1,b_2)$ with $b_0=0$ and $2b_2=\frac75b_1$. The remaining statements follow from (ii): two logarithmic solutions with the same $g$ differ by a power-series solution, and the normalisation of $h$ fixes it.
\end{proof}

\begin{defn}[Conifold frames]
\label{def:set-frames}
A \emph{conifold frame} is given by $(\alpha,\beta)\in\bbQ^2$:
\[
\varpi_0\coloneqq f_0+\alpha\varpi_1+\beta f_2,\qquad\tc\coloneqq\varpi_1/\varpi_0 .
\]
Then $\varpi_0=1+\alpha u+O(u^2)$, so $\varpi_0(0)=1$ and $w_1\coloneqq\coef{u^1}\varpi_0=\alpha$, and $\tc=u+(\frac7{10}-\alpha)u^2+O(u^3)$, so $\tc'(0)=1$, $\Dq\tc(0)=-1$ and $\gamma_t+w_1=\frac7{10}$ for $\gamma_t\coloneqq\coef{u^2}\tc$. The vanishing period is always normalised by $\varpi_1=u+O(u^2)$. In the frame $(\alpha,\beta)=(0,0)$, $\varpi_0=1+O(u^3)$ and $\tc=u+\frac7{10}u^2+\frac{41}{75}u^3+\frac9{20}u^4+O(u^5)$. These are \emph{the conifold frames} of the note. In them the normalisation (N0) of \cref{lem:gau-diag} holds; it uses only $\varpi_0(0)=1$, $\tc'(0)=1$, $\varpi_1=u+O(u^2)$, and $v(u)=u+O(u^2)$ (\cref{def:set-tau}).
\end{defn}

\begin{defn}[The differential ring of the periods and the conifold expansion]
\label{def:set-confexp}
Let
\[
\Nring\coloneqq\bbQ[L^{\pm1},\logq,c_0,c_1,c_2,c_3,i_0,i_1,i_2][c_0^{-1},\mathsf W^{-1}],\qquad\mathsf W\coloneqq c_0^2+c_0i_1-c_1i_0,
\]
the localisation of a polynomial ring. Its derivation $\Dq$ is fixed on the generators: $\Dq L\coloneqq(L^6-L)/5$, $\Dq\logq\coloneqq1$, $\Dq c_j\coloneqq c_{j+1}$ ($j\leq2$), $\Dq c_3\coloneqq(L^5-1)\bigl(2c_3+\frac75c_2+\frac25c_1+\frac{24}{625}c_0\bigr)$, which is $\PFfour c_0=0$ with $u=L^{-5}$, $\Dq i_j\coloneqq i_{j+1}$ ($j\leq1$), and $\Dq i_2\coloneqq i_3$, the element of $\Nring$ determined by the concomitant relation $\Conc(c_0,I_1)=0$, where $I_1\coloneqq c_0\logq+i_0$, $\Dq^jI_1=\logq\,c_j+jc_{j-1}+i_j$ for $j\leq3$, and $\Conc$ is the concomitant of $\PFfour$ (\cref{def:gau-conc}): the terms with $\logq$ cancel because $\Conc(c_0,c_0)=0$ identically, and the relation is linear in $i_3$ with the coefficient $-uc_0=-L^{-5}c_0$, a unit of $\Nring$. In $\Nring$ put $C_0\coloneqq c_0$, $I_1\coloneqq c_0\logq+i_0$, $T\coloneqq I_1/C_0$ and $C_1\coloneqq\Dq T=\mathsf W/c_0^2$. The \emph{evaluation} $\ev\colon\Nring\to\bbQ\fps q[\logq]$, $c_j\mapsto\Dq^jC_0$, $i_j\mapsto\Dq^j\tilde I_1$, $\logq\mapsto\logq$, $L\mapsto L$, is a homomorphism of differential rings, with $\ev(\mathsf W)=C_0^2C_1$: it respects $\Dq c_3$ because $\PFfour C_0=0$, and $\Dq i_2$ because $\Conc(C_0,I_1)=0$. Indeed $\PFfour C_0=\PFfour I_1=0$ (\cref{lem:set-periods}), so the identity of differential polynomials of \cref{lem:gau-conc}~(i), with \cref{lem:gau-adj}, gives $\Dq\Conc(C_0,I_1)=0$, and $\Conc(C_0,I_1)$ is a constant; its part of order $q^0$ vanishes, because there only the leading coefficient of $\PFfour$ survives and the jets of $C_0$ and $I_1$ reduce to $(1,0,0,0)$ and $(\logq,1,0,0)$.

The \emph{conifold expansion} in the frame $(\alpha,\beta)$ is the ring homomorphism $\Cexp_{\alpha,\beta}\colon\Nring\to\bbQ\lser s$ with
\[
c_j\mapsto\Dq^j\varpi_0,\qquad i_j\mapsto\Dq^j\varpi_1-j\,\Dq^{j-1}\varpi_0+\Xi\,\Dq^j\varpi_0,\qquad\logq\mapsto\log(1-u)=-\Xi,\qquad L\mapsto s^{-1},
\]
where $\Xi\coloneqq-\log(1-u)=\sum_{n\geq1}u^n/n$; it is well defined because $c_0\mapsto\varpi_0$ and $\mathsf W\mapsto\varpi_0\Dq\varpi_1-\varpi_1\Dq\varpi_0=\varpi_0^2\Dq\tc$ are units ($\varpi_0(0)=1$, $\Dq\tc(0)=-1$). It maps $C_0\mapsto\varpi_0$, $I_1\mapsto\varpi_1$, $T\mapsto\tc$ and $C_1\mapsto\Dq\tc$. No relation beyond $\PFfour$ and the concomitant relation is imposed. Replacing the primitive $-\Xi$ of $1$ by $-\Xi+k$, $k\in\bbQ$, composes $\Cexp_{\alpha,\beta}$ with the automorphism $i_j\mapsto i_j-kc_j$, $\logq\mapsto\logq+k$ of $\Nring$; this automorphism commutes with $\Dq$ and fixes $c_0$, $I_1$, $\mathsf W$ and $L$, hence $C_0$, $C_1$, $\Delta_i^{1/2}=\zeta^{4i}C_0/L$, the rows (\cref{cor:gau-rowsNN}) and the genus-$g$ sum $F_g^{\Nring}$ of \cref{def:mai-Fcf}. That $\Cexp_{\alpha,\beta}$ commutes with $\Dq$ is \cref{lem:gau-confexp}, and that the rows lie in $\Nring\otimes\bbQ(\zeta)$ is \cref{cor:gau-rowsNN}. An element of $\Nring\otimes\bbQ(\zeta)$, or a series with coefficients in it, is \emph{regular at the conifold point} if its conifold expansion lies in $\bbQ\fps u$, coefficientwise. The comparison with the conifold expansion of \cite{CGYZ26} is \cref{rmk:mai-constant}.
\end{defn}

\begin{defn}[The natural variable]
\label{def:set-tau}
Put
\[
\tau(u)\coloneqq\frac1{{}_2F_1(1,\frac45;\frac95;u)}=\Bigl(\sum_{n\geq0}\frac{4u^n}{5n+4}\Bigr)^{-1}\in1+u\bbQ\fps u,\qquad v(u)\coloneqq u\,\tau(u)^{-5/4}\in u+u^2\bbQ\fps u,
\]
the power being the binomial series, and $\varsigma\coloneqq v^{1/5}=s\,\tau^{-1/4}$, so that $z/\varsigma^4=\tau(u)\zh$. For instance $\tau=1-\frac49u+O(u^2)$ and $v=u+\frac59u^2+O(u^3)$. The letter $\tau$ denotes only this series.
\end{defn}

\begin{lem}[The canonical coordinates at the conifold point]
\label{lem:set-canon}
In $\bbQ(\zeta)\lser s$, with $\Dq=-(1-u)\,\dd/\dd u$ and $L=s^{-1}$, the solutions $y$ of $\Dq y=\zeta^iL$ are the series $y=c-\frac54\zeta^iu^{4/5}/\tau(u)$, $c\in\bbQ(\zeta)$, where $u^{4/5}=s^4$. So in the conifold chart the canonical coordinates of \cref{cor:set-cyclicPF5}~(a), which the note uses only through the equation $\Dq u_i=\zeta^iL$ and which are determined by it up to additive constants, are the series $u_i=u_i^{\mathrm c}-\frac54\zeta^iu^{4/5}/\tau(u)$, with constants $u_i^{\mathrm c}$, their values at $u=0$; and
\[
\tau(u)\,\zh=-\frac54\,\frac{\zeta^iz}{u_i-u_i^{\mathrm c}},\qquad u_i-u_i^{\mathrm c}=\zeta^i\bigl(u_0-u_0^{\mathrm c}\bigr).
\]
\end{lem}

\begin{proof}
Termwise, $\int_0^ut^{-1/5}(1-t)^{-1}\dd t=\sum_{n\geq0}\int_0^ut^{n-1/5}\dd t=\sum_n\frac{u^{n+4/5}}{n+4/5}=\frac54u^{4/5}\sum_n\frac{4u^n}{5n+4}=\frac54u^{4/5}/\tau(u)$. So $\frac{\dd}{\dd u}\bigl(\frac54u^{4/5}/\tau\bigr)=u^{-1/5}/(1-u)$, and $\Dq\bigl(-\frac54\zeta^iu^{4/5}/\tau\bigr)=\zeta^iu^{-1/5}=\zeta^iL$. The constants of $\Dq$ on $\bbQ(\zeta)\lser s$ are $\bbQ(\zeta)$, since $\Dq f=0$ means $\dd f/\dd s=0$; this gives the solutions. By \cref{cor:set-cyclicPF5}~(a), $\Dq u_i=\zeta^iL$. Finally $\zh=z/s^4=z/u^{4/5}$ gives $-\frac54\zeta^iz/(u_i-u_i^{\mathrm c})=\tau z/u^{4/5}=\tau\zh$.
\end{proof}

The constants $u_i^{\mathrm c}$ play no role: the note uses the canonical coordinates only through $\Dq u_i=\zeta^iL$ and through the differences $u_i-u_i^{\mathrm c}$, the integrals from the conifold point, whose orientation is fixed by $\Dq u_i=\zeta^iL$. The exponentials $\re^{u_i/z}$, and with them the constant factors $\re^{u_i^{\mathrm c}/z}$, are stripped from the formal solutions and are never absorbed into a normalisation series (\cref{thm:gau-R0}).

\input{sections/s02_setup_6}



